\section{Decisiones y consideraciones de diseño}

Las decisiones adoptadas para completar la información no especificada por
los mini--mundos son las siguientes:

\begin{itemize}
  \item En el ejercicio 2-i se distingue entre registrar una autorización
        general para manejar un taxi y registrar la asignación efectiva de un
        chofer y un taxi dentro de un sitio.
  \item En el modelo universitario se conserva la estructura original cuando
        el enunciado no solicita modificarla. Las flechas expresan unicidad y
        la participación se distingue únicamente mediante líneas simples o
        dobles, conforme a la notación empleada por el equipo.
  \item En los números racionales se adopta una especialización total y
        disjunta entre fracciones reducidas y no reducidas. El factor se
        conserva como atributo derivado de la relación \textbf{Reducir}.
  \item En el metamodelo E--R, \textbf{Corresponder} reúne las restricciones de
        cardinalidad y participación de cada intervención de una entidad en
        una relación. Los tipos de atributo se consideran parte de la
        descripción resumida por el atributo \emph{Atributos}.
  \item En el SIG se usan los identificadores indicados por el enunciado; las
        longitudes se interpretan en kilómetros y las alturas en metros. La
        especialización de montaña es total y disjunta, y la categoría
        \textbf{Elemento geológico} permite monitorear ríos o montañas sin
        duplicar la relación.
\end{itemize}

Los cuatro modelos solicitados se entregan en sus archivos originales, tanto
en formato editable de Draw.io como en PNG. Las figuras se incorporan también al documento para que
las restricciones y sus justificaciones puedan revisarse de manera conjunta.
